\section{Installation and test suite}
\label{sec:install}

\subsection{Build requirements}
\begin{itemize}
    \item CMake $\geq$ 3.18, C++17 compiler, FFTW3, OpenMP.
    \item Optional CUDA toolkit for \texttt{BUILD\_CUDA=ON} (Tesla T4, \texttt{sm\_75}).
    \item Python 3.10+ with \texttt{matplotlib} and \texttt{pandas} for figure regeneration.
\end{itemize}

\subsection{Build matrix}
\begin{verbatim}
cmake -S . -B build -DCMAKE_BUILD_TYPE=Release
cmake --build build -j
# Optional GPU:
cmake -S . -B build -DCMAKE_BUILD_TYPE=Release -DBUILD_CUDA=ON
\end{verbatim}

\subsection{Test suite}
\texttt{./build/pic\_test} must report:
\begin{verbatim}
PASS CIC charge conservation
PASS TSC charge conservation
PASS Esirkepov charge conservation
PASS NGP charge conservation
PASS GPU CIC atomics charge conservation      # CUDA builds only
PASS GPU CIC privatized charge conservation   # CUDA builds only
\end{verbatim}
Exit code 0 indicates success.

\subsection{Example workflows}
The \texttt{examples/} directory provides:
\begin{itemize}
    \item \texttt{examples/cic\_benchmark.sh}: minimal CIC microbenchmark reproducing one archived CSV row.
    \item \texttt{examples/run\_validation.sh}: Langmuir + conservation + two-stream validation.
\end{itemize}
After running benchmarks, archive CSVs to \texttt{data/benchmarks/} and execute \texttt{python3 scripts/plot\_results.py}.
